\documentclass[10pt,a4paper]{amsart}
\usepackage[margin=19mm]{geometry}
\usepackage{amsmath,amssymb,mathtools}
\usepackage[scr=rsfs,cal=euler]{mathalfa}
\usepackage{xfrac}
\usepackage[unicode,hidelinks,bookmarks=false]{hyperref}
\newcommand{\ie}{i.e.\ }
\newcommand{\cf}{cf.\ }
\newcommand{\resp}{respectively\ }
\newcommand{\Subsection}{\subsection}
\providecommand{\nobreakdash}{\nobreak\hbox{-}}
\setlength{\emergencystretch}{2em}
\multlinegap0pt
\usepackage{a4wide}
\usepackage{lmodern}
\usepackage{mathtools,enumitem,url}
\usepackage{flafter,float}
\usepackage{tikz}
\usetikzlibrary{arrows.meta,positioning,calc}
\setlength{\emergencystretch}{3em}
\newtheorem{theorem}{Theorem}[section]
\newtheorem{proposition}[theorem]{Proposition}
\newtheorem{lemma}[theorem]{Lemma}
\newtheorem{corollary}[theorem]{Corollary}
\newtheorem*{theoremA}{Theorem A}
\newtheorem*{theoremB}{Theorem B}
\newtheorem{definition}[theorem]{Definition}
\newtheorem{remark}[theorem]{Remark}
\newcommand{\R}{\mathbb R}
\newcommand{\N}{\mathbb N}
\newcommand{\supp}{\operatorname{supp}}
\newcommand{\Span}{\operatorname{span}}
\newcommand{\one}{\mathbf 1}
\newcommand{\eps}{\varepsilon}
\newcommand{\restr}[1]{{\restriction_{#1}}}
\numberwithin{equation}{section}
\begin{document}
\title[A consistent failure of separable quotients for pointwise fu]{A consistent failure of separable quotients for pointwise function spaces}
\author{Tomasz Kania; Jerzy Kąkol}
\date{}
\maketitle
\noindent\emph{Source and licence.} A consistent failure of separable quotients for pointwise function spaces. Version arXiv:2609.19351v1 (CC BY 4.0). This material is distributed under the Creative Commons Attribution 4.0 International licence, http://creativecommons.org/licenses/by/4.0/, as recorded in the arXiv OAI licence metadata for this exact version and shown on the abstract page https://arxiv.org/abs/2609.19351v1. Full rights notice: licenses/arxiv-2609.19351-CC-BY-4.0.txt.\par\smallskip
\noindent\textit{Editorial scope.} Counted unit: the original \texttt{lemma} environment \texttt{lem:support-projection} at source lines 427--441 together with its complete proof at lines 442--456; the delivered document keeps source lines 415--457 so that every referenced label and every citation resolves inside it. Native editable source is the paper's own concatenated TeX; the AMS wrapper is editorial formatting only. No publisher class, upstream script or unrelated artwork is redistributed.
\section{Prediction and preservation at countable limits}\label{sec:limits}

We prepare two ingredients for the construction of $K$. First, diamond
will predict every countable sequence of finitely supported measures
through its coordinate projections. Secondly, we specify the convergence
requirements to be maintained throughout the inverse system and show
that they extend to each countable limit stage.

\subsection{Prediction of measure sequences}\label{subsec:prediction}

We shall apply diamond to sequences of finitely supported measures on $2^{\omega_1}$.  The union of the supports of one such sequence is countable.  The next observation ensures that sufficiently long coordinate projections preserve all its atoms, coefficients and variation norms.


\begin{lemma}\label{lem:support-projection}
Let $S\subseteq2^{\omega_1}$ be countable, and let
$p_\alpha:2^{\omega_1}\to2^\alpha$ be restriction of coordinates.
There is $\eta<\omega_1$ such that, whenever
$\eta\leqslant\alpha<\omega_1$, the map $p_\alpha$ is injective on $S$.
For every such $\alpha$, the map
\[
 (p_\alpha)_*:L(S)\longrightarrow L(p_\alpha[S]),
 \qquad \sum_{j<r}a_j\delta_{x_j}
       \longmapsto\sum_{j<r}a_j\delta_{p_\alpha(x_j)},
\]
is a linear isometric bijection, where $L(S)$ denotes the finitely
supported measures whose supports are contained in $S$; no closedness
of $S$ is required.
\end{lemma}

\begin{proof}
For every pair of distinct points $x,y\in S$, choose
$\xi(x,y)<\omega_1$ with $x(\xi(x,y))\neq y(\xi(x,y))$.
The set of these coordinates is countable.  Choose a countable
$\eta\geqslant\omega$ strictly greater than every chosen coordinate.
Then $p_\alpha(x)\neq p_\alpha(y)$ whenever $\alpha\geqslant\eta$.
Consequently, if $\mu=\sum_{j<r}a_j\delta_{x_j}$ is a reduced finite
representation with $x_j\in S$, then its push-forward has the same
nonzero coefficients at distinct points, and hence
\[
 \|(p_\alpha)_*\mu\|=\sum_{j<r}|a_j|=\|\mu\|.
\]
Every finitely supported measure on $p_\alpha[S]$ has a unique such
lift to $S$, proving the last assertion.
\end{proof}

\end{document}
